\makeabstract
{
Making Intervention Do-ABLE: Analysis of Police Officers’ Perceptions of Active Bystandership Training
}
{
Jiaqi Huang (1), Viktoria E. Leopold (1), Cassandra Ramdath (2), Catherine A. Sanderson (1,2)
}
{
\textsuperscript{1}Department of Psychology, Amherst College, Amherst, MA, USA\\
\textsuperscript{2}Active Bystandership for Law Enforcement (ABLE) Program, Georgetown Law, Washington, DC, USA
}
{
Purpose: 
This study investigates large-scale officer feedback on what agencies, direct supervisors, and each individual officer should do to improve active bystandership. This will serve to fix the research gap on how officers themselves perceive the effectiveness of ABLE training, and what they believe can strengthen its impact.
}
{
Methods:
This qualitative study analyzed responses from 48,077 officers across over 400 police departments in the USA and Canada to an open-ended national survey that explored perceptions, attitudes, and beliefs about active bystandership and peer intervention. The survey was administered following an 8-hour virtual officer training program. Researchers analyzed each question thematically based on an established codebook, tested for interrater reliability, and used it to identify patterns and themes related to how individual officers, supervisors, and agencies can support active bystandership.
}
{
Results:
To promote active bystandership, officers stressed ongoing training along with real-world application, and though most found ABLE effective or already in use, some suggested shorter, in-person training. They emphasized open discussion across all ranks, reflections after incidents, regular reminders to reinforce intervention, and outreach via promotions to share ABLE with other communities. Officers urged leadership to verbally support active bystandership through exemplary behavior, direct field involvement, and called for accountability. They also asked for recognition of successful interventions through formal or verbal awards. Officers called for mandatory ABLE training and policies with firm anti-retaliation protections, expressing that confidence in intervention increases when officers are shielded from backlash and negative consequences. They stressed the importance of individual action through speaking up and intervening, openness to feedback across all ranks, and continuous personal growth to build confidence and intervention skills. Officers want a culture of awareness and supportive relationships that promote intervention while addressing mental and physical well-being challenges created by toxic leadership, staffing shortages, and negative work environments.
}
{
Conclusion:
Departments can use this information to refine ABLE’s format, while also implementing the policies that support intervention, reducing police violence and misconduct. Promoting ABLE and sharing successful stories also fosters accountability to rebuild public trust in law enforcement. Future research could examine connections between demographic data and officers’ responses and recommendations for encouraging ABLE. 
}

\newpage
